% =====================================================================
%  PRÉAMBULE
%  Contraintes ENSIM : Times New Roman 12, interligne simple,
%  marges 2 cm, pagination, labels figures/tableaux >= 9 pt.
% =====================================================================

% ---- Encodage et langue ----
\usepackage[utf8]{inputenc}
\usepackage[T1]{fontenc}
\usepackage[french]{babel}
\usepackage{csquotes}

% ---- Police : équivalent Times New Roman ----
\usepackage{newtxtext}
\usepackage{newtxmath}

% ---- Mise en page : marges 2 cm ----
\usepackage[a4paper,margin=2cm,headheight=15pt]{geometry}

% ---- Interligne simple (valeur par défaut) ----
% Pour un interligne légèrement aéré, décommenter :
% \usepackage{setspace}\setstretch{1.15}

% ---- Graphiques et couleurs ----
\usepackage{graphicx}
\graphicspath{{images/}}
\usepackage[dvipsnames,table]{xcolor}
\usepackage{float}
% ---- Diagrammes (schémas de structure, sans images externes) ----
\usepackage{tikz}
\usetikzlibrary{shapes.geometric,arrows.meta,positioning,fit,backgrounds,shadows}

% ---- Palette VINCI Energies (bleu / rouge) ----
\definecolor{vinciblue}{HTML}{003366}
\definecolor{vincired}{HTML}{E2001A}
\definecolor{grisfonce}{HTML}{333333}
\definecolor{grisclair}{HTML}{F2F2F2}
\definecolor{azurblue}{HTML}{0078D4}
\definecolor{phrec}{HTML}{B45309}
\definecolor{phprod}{HTML}{15803D}
% Commande utilitaire pour les sous-rôles dans les schémas TikZ
\newcommand{\rrole}[1]{\scriptsize\textcolor{grisfonce}{#1}}

% ---- Tableaux ----
\usepackage{array}
\usepackage{booktabs}
\usepackage{tabularx}
\usepackage{multirow}
\usepackage{longtable}
\renewcommand{\arraystretch}{1.3}

% ---- Légendes (labels >= 9 pt, gras sur le numéro) ----
\usepackage[font=small,labelfont=bf,labelsep=colon]{caption}
\usepackage{subcaption}

% ---- En-têtes et pieds de page ----
\usepackage{fancyhdr}
% Pages blanches d'insertion totalement vides (verso blanc)
\usepackage{emptypage}
\pagestyle{fancy}
\fancyhf{}
\renewcommand{\headrulewidth}{0.4pt}
\fancyhead[L]{\footnotesize\itshape\leftmark}
\fancyhead[R]{\footnotesize ENSIM \textbar{} VESI}
\fancyfoot[C]{\thepage}
\fancypagestyle{plain}{%
  \fancyhf{}%
  \fancyfoot[C]{\thepage}%
  \renewcommand{\headrulewidth}{0pt}%
}

% ---- Titres de chapitres et sections (sobres, sans traits) ----
\usepackage{titlesec}
\titleformat{\chapter}[display]
  {\normalfont\LARGE\bfseries\color{vinciblue}}
  {\filright\large\chaptertitlename\ \thechapter}
  {1ex}
  {\filright}
\titlespacing*{\chapter}{0pt}{10pt}{20pt}

\titleformat{\section}
  {\normalfont\normalsize\bfseries\color{vinciblue}}{\thesection}{1em}{}
\titleformat{\subsection}
  {\normalfont\small\bfseries\color{grisfonce}}{\thesubsection}{1em}{}

% ---- Listes plus compactes ----
\usepackage{enumitem}
\setlist{itemsep=2pt,topsep=3pt}

% ---- Extraits de code (listings, robuste sur Overleaf) ----
\usepackage{listings}
\lstdefinestyle{codestyle}{
  basicstyle=\ttfamily\footnotesize,
  keywordstyle=\color{vinciblue}\bfseries,
  commentstyle=\color{gray}\itshape,
  stringstyle=\color{vincired},
  numberstyle=\tiny\color{gray},
  numbers=left,
  numbersep=8pt,
  breaklines=true,
  showstringspaces=false,
  frame=single,
  rulecolor=\color{gray!40},
  backgroundcolor=\color{grisclair},
  tabsize=2,
  captionpos=b,
}
\lstset{style=codestyle}
% C# : \lstset{language=[Sharp]C}
% SQL : \lstset{language=SQL}

% ---- Bibliographie (style numérique [1], [2]...) ----
\usepackage[backend=biber,style=numeric,sorting=none]{biblatex}

% ---- Liens hypertexte (chargé tard) ----
\usepackage[hidelinks]{hyperref}
\hypersetup{
  colorlinks=true,
  linkcolor=vinciblue,
  citecolor=vincired,
  urlcolor=vinciblue,
  pdftitle={Rapport de PFE - Ahmed BOURAZZA},
  pdfauthor={Ahmed BOURAZZA},
}
\usepackage[nameinlink,french]{cleveref}

% Table des matières / figures / tableaux et renvois internes en noir
% (les citations bibliographiques [1] restent en rouge)
\AtBeginDocument{\hypersetup{linkcolor=black}}

% ---- Microtypographie ----
\usepackage{microtype}

% ---- Commandes utiles ----
\usepackage{xspace}
\newcommand{\csbag}{CS~BAG\xspace}

% Capture d'écran : affiche un cadre indicatif tant que le fichier est absent.
\newcommand{\screenshot}[2][0.85\textwidth]{%
  \IfFileExists{images/#2}%
    {\includegraphics[width=#1]{#2}}%
    {\fbox{\parbox[c][0.25\textheight][c]{\dimexpr#1-2\fboxsep-2\fboxrule\relax}%
      {\centering\small\itshape Capture à insérer\\[3pt]\texttt{images/#2}}}}%
}